有了密度或分数之后，还不能直接说已经获得了样本。
MCMC 的想法是构造一个容易执行的转移核 $K(x,\mathrm dy)$，
反复从当前状态走到下一状态，使目标分布 $p$ 成为这条链的不变分布：
\[
p(y)=\int p(x)K(x,y)\,\mathrm dx.
\]
如果初始状态不是从 $p$ 抽出的，还需要链逐渐接近这个不变分布。
因此，“$p$ 是不变分布”与“有限步结果已经服从 $p$”是两件事。

与分数模型最直接相关的是 Langevin 动力学。先写连续过程
\[
\mathrm dX_r=\tfrac12\nabla_x\log p(X_r)\,\mathrm dr+\mathrm dW_r.
\]
这里的 $r$ 是采样过程的时间，还不是后面扩散加噪使用的噪声时间。
将它的漂移与扩散系数代入后面要推导的 Fokker--Planck 方程，
若令当前密度为 $p$，右边为
\[
-\tfrac12\nabla\cdot(p\nabla\log p)+\tfrac12\Delta p=0,
\]
所以 $p$ 确实是不变密度。再在满足遍历性等条件时讨论趋近目标分布。

以步长 $h>0$ 做离散近似，得到
\[
x_{k+1}=x_k+\tfrac h2 s(x_k)+\sqrt h\,z_k,
\qquad z_k\sim\mathcal N(0,I).
\]
分数项将样本推向对数密度较高的位置，随机项使样本能够探索不同区域。
只沿分数上升会不断靠近局部模式，不能替代按整个概率分布采样。
固定步长还会产生离散化偏差；分布有相隔很远的多个模式时，混合也可能很慢。
NCSN 后面采用不同噪声级的退火采样，正是继续处理这些困难。
